% =============================================================================
% exercises/ejercicios_cap10.tex
% Ejercicios propuestos — Capítulo 10: Química en la Troposfera
% =============================================================================
\newpage
\section*{Ejercicios Propuestos}
\addcontentsline{toc}{section}{Ejercicios propuestos}
\label{sec:ejercicios_cap10}

\begin{bloque_ejercicios}[Ejercicios --- Cap\'itulo 10: Qu\'imica Troposf\'erica]

\begin{ejercicios}

% ----- Ejercicio 1: Producción del radical OH ------------------------------
\item La producci\'on del radical OH en la troposfera sigue la
secuencia de tres reacciones (\ref{OHR1}--\ref{OHR3}):
\[
    \ce{O3 ->[$h\nu$] O2 + O(^1D)} \quad
    \ce{O(^1D) + M -> O(^3P) + M} \quad
    \ce{O(^1D) + H2O -> 2OH}
\]

\begin{enumerate}[label=\alph*)]
    \item La producci\'on simplificada de OH es:
          \[
              P_{\mathrm{OH}} \approx
              \frac{2k_1 k_3}{k_2[\mathrm{M}]}[\ce{O3}][\ce{H2O}]
          \]
          Explique por qu\'e la aproximaci\'on $k_2[\mathrm{M}] \gg
          k_3[\ce{H2O}]$ es v\'alida en la troposfera.

    \item La producci\'on de \ce{O(^1D)} s\'olo ocurre con radiaci\'on
          entre \SIrange{300}{320}{\nano\metre}. Explique por qu\'e
          la radiaci\'on de $\lambda < \SI{300}{\nano\metre}$ no
          alcanza la troposfera y por qu\'e $\lambda > \SI{320}{\nano\metre}$
          no produce \ce{O(^1D)}.

    \item La concentraci\'on media global de OH es
          $\overline{[\mathrm{OH}]} = 1.2\times10^6$~mol\'ec/cm$^3$.
          Usando $\tau_i = 1/(k_i\overline{[\mathrm{OH}]})$,
          calcule el tiempo de vida del CO
          ($k_{\mathrm{CO}} = 2.4\times10^{-13}$~cm$^3$~mol\'ec$^{-1}$~s$^{-1}$)
          y del \ce{CH4}
          ($k_{\mathrm{CH_4}} = 6.4\times10^{-15}$~cm$^3$~mol\'ec$^{-1}$~s$^{-1}$).
          Exprese los resultados en d\'ias y a\~nos.

    \item Durante la campa\~na MCMA-2003 en la Ciudad de M\'exico
          se midieron picos de $[\mathrm{OH}] = 180\times10^6$~mol\'ec/cm$^3$.
          Calcule el tiempo de vida del CO en esas condiciones.
          ?`C\'omo cambia respecto al valor global?
\end{enumerate}

\textit{[Respuesta:
(c) $\tau_{\mathrm{CO}} = 1/(2.4\times10^{-13}
    \times 1.2\times10^6) = 3.47\times10^6$~s $\approx$ 40~d\'ias;
    $\tau_{\mathrm{CH_4}} = 1/(6.4\times10^{-15}
    \times 1.2\times10^6) = 1.3\times10^8$~s $\approx$ 4.1~a\~nos;
(d) $\tau_{\mathrm{CO}} = 1/(2.4\times10^{-13}
    \times 1.8\times10^8) = 23$~s --- reducido dram\'aticamente]}

% ----- Ejercicio 2: Ciclo HO_x y producción de ozono (oxidación del CO) ---
\item El mecanismo de oxidaci\'on del CO en presencia de \ce{NO_$x$}
produce \ce{O3} de forma neta. Considere la secuencia de reacciones
del cap\'itulo:
\begin{eqnarray*}
    \ce{CO + OH -> CO2 + H} &
    \ce{H + O2 ->[M] HO2} \\
     \ce{O^. + O2 + M -> O3 + M}&
    \ce{NO2 ->[$h\nu$] NO + O^.}\\
    \ce{HO2 + NO -> OH + NO2}
\end{eqnarray*}

\begin{enumerate}[label=\alph*)]
    \item Sume las cinco reacciones para obtener la reacci\'on neta
          global. Verifique que coincide con:
          \ce{CO + 2O2 -> CO2 + O3}.

    \item Identifique los cat\'alizadores del ciclo (especies
          que se regeneran) y los reactivos que se consumen neta\-mente.

    \item Si $[\mathrm{NO}] = 1\times10^{11}$~mol\'ec/cm$^3$
          y $k_{\ce{HO2+NO}} = 8.1\times10^{-12}$~cm$^3$~mol\'ec$^{-1}$~s$^{-1}$,
          calcule la velocidad de regeneraci\'on de OH (en
          mol\'ec~cm$^{-3}$~s$^{-1}$) suponiendo
          $[\ce{HO2}] = 5\times10^8$~mol\'ec/cm$^3$.

    \item Explique por qu\'e la reacci\'on alternativa
          \ce{HO2 + HO2 -> H2O2 + O2} termina la cadena y
          reduce la capacidad oxidativa de la atm\'osfera.
          ?`En qu\'e condiciones domina este camino frente
          a la ruta con NO?
\end{enumerate}

\textit{[Respuesta:
(c) $v = 8.1\times10^{-12}
    \times 1\times10^{11}\times5\times10^8
    = 4.05\times10^8$~mol\'ec~cm$^{-3}$~s$^{-1}$;
(d) domina cuando $[\mathrm{NO}]$ es bajo
    (atm\'osfera limpia o remota)]}

% ----- Ejercicio 3: Oxidación del metano — estados de oxidación del C -----
\item La oxidaci\'on del \ce{CH4} produce una secuencia de
intermediarios que elevan el estado de oxidaci\'on del carbono
de \ce{C^{-IV}} hasta \ce{C^{+IV}}.

\begin{enumerate}[label=\alph*)]
    \item Complete la tabla de intermediarios y estados de oxidaci\'on
          del carbono en la cadena \ce{CH4 -> CH3OOH -> CH2O -> CO -> CO2}:

{%
\centering
\begin{tabular}{lcl}
\toprule
\textbf{Compuesto} & \textbf{Estado ox. de C} &
\textbf{Reacci\'on que lo produce} \\
\midrule
\ce{CH4}    & $-$IV & emisi\'on biol\'ogica \\
\ce{CH3OOH} &       & \\
\ce{CH2O}   &       & \\
\ce{CO}     &       & \\
\ce{CO2}    & $+$IV & \\
\bottomrule
\end{tabular}
\par
}%

    \item En atm\'osfera con alto \ce{NO_$x$}, la oxidaci\'on neta es:
          \ce{CH4 + 10O2 -> CO2 + H2O + 5O3 + 2OH}.
          ?`Cu\'antas mol\'eculas de \ce{O3} se producen por cada
          mol\'ecula de \ce{CH4} oxidada? ?`Y cu\'antas de OH?

    \item En atm\'osfera con bajo \ce{NO_$x$}, la oxidaci\'on neta es:
          \ce{CH4 + 3OH + 2O2 -> CO2 + 3H2O + HO2}.
          Compare los balances de OH en ambas condiciones.
          ?`En cu\'al se consume m\'as OH?

    \item Explique en t\'erminos mecanicistas por qu\'e el \ce{NO_$x$}
          es cr\'itico para mantener las concentraciones de \ce{O3}
          y OH en la troposfera.
\end{enumerate}

\textit{[Respuesta:
(a) \ce{CH3OOH}: $-$II; \ce{CH2O}: 0; \ce{CO}: $+$II; \ce{CO2}: $+$IV;
(b) 5 mol\'eculas \ce{O3} y 2 mol\'eculas OH;
(c) bajo NOx: se consumen 3 OH netos; alto NOx: se generan 2 OH]}

% ----- Ejercicio 4: Fuentes de NOx y ciclo diurno -------------------------
\item Con base en el \textbf{Cuadro~\ref{noxt}} de emisiones de \ce{NO_$x$}:

\begin{enumerate}[label=\alph*)]
    \item Calcule el total de emisiones de \ce{NO_$x$} en
          \si{\tera\gram}/a\~no sumando todas las fuentes.
          ?`Qu\'e porcentaje corresponde a fuentes antropog\'enicas
          (combusti\'on f\'osil + aviaci\'on)?

    \item El ciclo diurno del \ce{NO_$x$} involucra:
          \ce{NO + O3 -> NO2 + O2} (diurno y nocturno) y
          \ce{NO2 ->[$h\nu$] NO + O^.} (solo diurno).
          Explique por qu\'e de noche todo el \ce{NO_$x$}
          se presenta como \ce{NO2}.

    \item El sumidero principal diurno del \ce{NO_$x$} es
          \ce{NO2 + OH ->[M] HNO3}.
          Si $k = 1.2\times10^{-11}$~cm$^3$~mol\'ec$^{-1}$~s$^{-1}$
          y $[\mathrm{OH}] = 3\times10^6$~mol\'ec/cm$^3$,
          calcule el tiempo de vida del \ce{NO2} frente
          a esta reacci\'on.

    \item El sumidero nocturno sigue la secuencia:
          \ce{NO2 + O3 -> NO3 + O2},
          \ce{NO3 + NO2 ->[M] N2O5},
          \ce{N2O5 + H2O -> 2HNO3}.
          Explique por qu\'e el \ce{HNO3} es un
          \textbf{sumidero terminal} del \ce{NO_$x$} en
          la troposfera (a diferencia de la estratosfera).
\end{enumerate}

\textit{[Respuesta:
(a) Total = 45.6~Tg/a\~no; antropog\'enico = 21.5~Tg
    (47\%);
(c) $\tau = 1/(1.2\times10^{-11}\times3\times10^6)
    = 2.78\times10^4$~s $\approx$ 7.7~h]}

% ----- Ejercicio 5: PAN como reservorio y transporte de NOx ----------------
\item El PAN (\ce{CH3C(O)OONO2}) es un reservorio de \ce{NO_$x$}
cuyo tiempo de vida depende fuertemente de la temperatura.

\begin{enumerate}[label=\alph*)]
    \item Escriba la secuencia de tres reacciones que forman PAN
          a partir del acetaldeh\'ido (\ce{CH3CHO}), incluyendo
          los intermediarios \ce{CH3CO^.} y \ce{CH3C(O)OO^.}
          (reacciones \ref{NOR34}--\ref{NOR36} del cap\'itulo).

    \item La descomposici\'on t\'ermica del PAN sigue
          cin\'etica de primer orden con $k = A\,e^{-E_a/RT}$.
          En el cap\'itulo se indica que el tiempo de vida es
          1~h a \SI{295}{\kelvin} y varios meses a \SI{250}{\kelvin}.
          Calcule la relaci\'on $k_{295}/k_{250}$ sabiendo que
          $\tau = 1/k$.

    \item Con base en la dependencia t\'ermica del PAN,
          explique por qu\'e el PAN puede transportar
          \ce{NO_$x$} a largas distancias cuando asciende
          a la troposfera media y superior.

    \item ?`Por qu\'e el PAN es un contaminante secundario
          m\'as relevante que el \ce{HNO3} para el transporte
          intercontinental de \ce{NO_$x$}? Mencione las propiedades
          f\'isicoquimicas que lo diferencian.
\end{enumerate}

\textit{[Respuesta:
(b) $\tau_{295} = 1$~h $= 3600$~s $\Rightarrow k_{295}
    = 2.78\times10^{-4}$~s$^{-1}$;
    $\tau_{250} \approx$ 3~meses $= 7.8\times10^6$~s
    $\Rightarrow k_{250} = 1.28\times10^{-7}$~s$^{-1}$;
    $k_{295}/k_{250} \approx 2\,200$ veces m\'as r\'apido a 295~K]}

% ----- Ejercicio 6: Radical NO3 — química nocturna ------------------------
\item El radical \ce{NO3} domina la qu\'imica oxidante nocturna.

\begin{enumerate}[label=\alph*)]
    \item Explique por qu\'e la concentraci\'on de \ce{NO3} es
          pr\'acticamente cero durante el d\'ia pero puede ser
          significativa por la noche. Mencione las dos razones
          principales del cap\'itulo.

    \item La reacci\'on \ce{NO3 + NO -> 2NO2} es una
          transferencia de \'atomo de O. Calcule el cambio en
          el estado de oxidaci\'on del N en el NO (reactivo)
          y en el \ce{NO2} (producto).

    \item Indique los productos principales de la reacci\'on
          \ce{NO3 + C2H4} seg\'un el cap\'itulo (ep\'oxido vs
          dinitrato) y explique cu\'al de los dos caminos
          requiere la participaci\'on de \ce{O2}.

    \item La reacci\'on \ce{CH3CHO + NO3 -> CH3CO^. + HNO3}
          es la abstracci\'on del H aldeh\'idico. Compare esta
          reacci\'on con la de \ce{CH3CHO + OH}. ?`Cu\'al
          produce \ce{HNO3} directamente? ?`Qu\'e sucede
          con el radical \ce{CH3CO^.} en ambos casos?
\end{enumerate}

\textit{[Respuesta:
(b) en NO: N tiene estado ox. $+2$; en \ce{NO2}: $+4$;
    se oxida en 2 unidades;
(d) con \ce{NO3}: produce \ce{HNO3} directamente;
    con OH: produce \ce{H2O}; en ambos casos \ce{CH3CO^.}
    reacciona con \ce{O2} para formar \ce{CH3C(O)OO^.}]}

% ----- Ejercicio 7: Mecanismo de oxidación de alcanos (n-butano) -----------
\item La oxidaci\'on del n-butano (\ce{n-C4H10}) con OH ilustra
el mecanismo general para alcanos.

\begin{enumerate}[label=\alph*)]
    \item El primer paso es la abstracci\'on de H:
          \ce{n-C4H10 + OH -> sec-C4H9^. + H2O} (85\%) o
          \ce{n-C4H10 + OH -> 1-C4H9^. + H2O} (15\%).
          Explique por qu\'e el carbono secundario es
          significativamente m\'as reactivo que el primario
          en t\'erminos de energ\'ia de enlace C--H.

    \item Escriba la secuencia de reacciones desde el radical
          2-butilo (\ce{sec-C4H9^.}) hasta los productos
          finales (MEK = metil-etil-cetona y \ce{HO2}),
          pasando por:
          (i) adici\'on de \ce{O2} para formar radical peroxi,
          (ii) reacci\'on con NO para formar radical alcoxi y \ce{NO2},
          (iii) reacci\'on con \ce{O2} del radical alcoxi.

    \item En el proceso anterior, ?`se regenera OH? ?`Por qu\'e
          es importante que la reacci\'on \ce{HO2 + NO -> OH + NO2}
          ocurra para completar el ciclo \ce{HO_$x$}?

    \item Explique la diferencia entre la v\'ia que produce
          nitrato de alquilo (\ce{RO2 + NO -> RONO2}) y
          la que produce radical alcoxi. ?`Cu\'al act\'ua
          como reacci\'on de \textbf{terminaci\'on} de la cadena?
\end{enumerate}

% ----- Ejercicio 8: Isopreno — hidrocarburo biogénico clave ----------------
\item El isopreno (\ce{CH2=C(CH3)-CH=CH2}) es el hidrocarburo
biol\'ogico m\'as emitido (50\% de las emisiones globales de NMVOC).

\begin{enumerate}[label=\alph*)]
    \item Las constantes de velocidad del isopreno con OH,
          \ce{O3} y \ce{NO3} a \SI{298}{\kelvin} son
          $k_{\ce{OH}} = 1.0\times10^{-10}$,
          $k_{\ce{O3}} = 1.3\times10^{-17}$ y
          $k_{\ce{NO3}} = 6.8\times10^{-13}$~cm$^3$~mol\'ec$^{-1}$~s$^{-1}$.
          Usando concentraciones t\'ipicas diurnas
          $[\ce{OH}] = 2\times10^6$,
          $[\ce{O3}] = 7\times10^{11}$,
          $[\ce{NO3}] \approx 0$ mol\'ec/cm$^3$,
          calcule el tiempo de vida del isopreno frente
          a cada oxidante y determine cu\'al domina durante el d\'ia.

    \item Calcule el tiempo de vida del isopreno frente a
          \ce{NO3} para condiciones nocturnas t\'ipicas
          $[\ce{NO3}] = 5\times10^8$~mol\'ec/cm$^3$
          y compare con el tiempo de vida diurno frente a OH.

    \item Los productos principales de la reacci\'on con OH
          son MACR (metacrole\'ina) y MVK (metil vinil cetona).
          Consulte la \textbf{Figura~\ref{isopOH}} e identifique
          qu\'e tipo de reacci\'on inicial (adici\'on) permite
          que el isopreno reaccione con OH en cuatro posiciones
          diferentes.

    \item La oxidaci\'on del isopreno con \ce{O3} genera
          radicales OH como subproducto. Explique la
          importancia de esta fuente \textbf{nocturna} de OH
          en regiones forestales.
\end{enumerate}

\textit{[Respuesta:
(a) $\tau_{\ce{OH}} = 1/(10^{-10}\times2\times10^6)
    = 5000$~s $\approx$ 1.4~h;
    $\tau_{\ce{O3}} = 1/(1.3\times10^{-17}\times7\times10^{11})
    = 1.1\times10^5$~s $\approx$ 30~h; OH domina diurno;
(b) $\tau_{\ce{NO3}} = 1/(6.8\times10^{-13}
    \times5\times10^8) = 2941$~s $\approx$ 49~min; nocturno tambi\'en r\'apido]}

% ----- Ejercicio 9: PANs — formación y toxicidad ---------------------------
\item Los peroxiacilnitratos (PANs) son contaminantes secundarios
formados por la oxidaci\'on de aldeh\'{\i}dos en presencia de
\ce{NO_$x$} (reacciones \ref{NOR34}--\ref{NOR36}).

\begin{enumerate}[label=\alph*)]
    \item Escriba las reacciones para la formaci\'on del PAN
          m\'as simple (\ce{CH3C(O)OONO2}) a partir del
          acetaldeh\'ido, identificando cada intermediario
          reactivo.

    \item El mecanismo incluye la reacci\'on paralela:
          \ce{CH3C(O)OO^. + NO -> CH3C(O)O^. + NO2}
          seguida de
          \ce{CH3C(O)O^. -> CH3^. + CO2}.
          Compare este camino con la formaci\'on de PAN:
          ?`cu\'al consume m\'as \ce{NO_$x$} y cu\'al lo
          \textbf{secuestra} temporalmente?

    \item Los PANs se denominan colectivamente
          \ce{RC(O)OONO2}. Identifique el grupo funcional
          que lleva el \'atomo de C central (carbonilo peroxiacilo)
          y explique por qu\'e este grupo hace a los PANs
          potentes fitot\'oxicos (t\'oxicos para plantas).

    \item Dado que el PAN se descompone t\'ermicamente,
          ?`en qu\'e condiciones ambientales (hora del d\'ia,
          estaci\'on, altitud) esperaria usted encontrar
          las mayores concentraciones de PAN? Justifique
          con la cin\'etica de la reacci\'on~\ref{NOR37}.
\end{enumerate}

% ----- Ejercicio 10 (avanzado): Estimación del balance global de NOx ------
\item \textbf{[Avanzado]} El \textbf{Cuadro~\ref{noxt}} lista
las fuentes globales de \ce{NO_$x$}. Con una vida media del
\ce{NO_$x$} de $\approx$~1~d\'ia y las emisiones totales
del cuadro:

\begin{enumerate}[label=\alph*)]
    \item Calcule la masa total de \ce{NO_$x$} en la atm\'osfera
          usando $M = F\cdot\tau$ (donde $F$ es la tasa de
          emisi\'on y $\tau$ el tiempo de residencia).
          Exprese el resultado en \si{\tera\gram}.

    \item La conversi\'on de \ce{NO_$x$} a \ce{HNO3} es el
          sumidero principal. Si la reacci\'on diurna es
          \ce{NO2 + OH ->[M] HNO3} con $k = 1.2\times10^{-11}$
          cm$^3$~mol\'ec$^{-1}$~s$^{-1}$ y
          $[\mathrm{OH}]_{\text{dia}} = 3\times10^6$~mol\'ec/cm$^3$
          (activo 12~h/d\'ia), calcule el tiempo de vida
          diurno del \ce{NO_$x$} frente a esta reacci\'on.

    \item Considerando que la producci\'on de ozono por
          \ce{NO_$x$} es de $\approx$~5~mol de \ce{O3} por
          mol de \ce{CH4} oxidado (r\'egimen de alto NOx),
          estime la producci\'on global de \ce{O3} troposf\'erico
          debida s\'olo a la oxidaci\'on del \ce{CH4},
          asumiendo que la emisi\'on global de \ce{CH4} es
          $\approx$ 550~Tg/a\~no.
          (Masa molar: \ce{CH4} = 16; \ce{O3} = 48~g/mol)

    \item Discuta por qu\'e la reducci\'on de emisiones de
          \ce{NO_$x$} puede tener efectos \textbf{no lineales}
          sobre las concentraciones de \ce{O3} dependiendo
          del r\'egimen de NOx/COV en una cuenca urbana.
\end{enumerate}

\textit{[Respuesta:
(a) Total emisi\'on = 45.6~Tg/a\~no;
    $\tau = 1$~d\'ia $= 2.74\times10^{-3}$~a\~nos;
    $M = 45.6\times2.74\times10^{-3}
    \approx 0.125$~Tg en equilibrio;
(b) $\tau = 1/(1.2\times10^{-11}\times3\times10^6
    \times0.5) = 5.6\times10^4$~s $\approx$ 15.5~h;
(c) moles \ce{CH4} = $550\times10^{15}/16
    = 3.44\times10^{16}$~mol;
    \ce{O3} = $5\times3.44\times10^{16}\times48\times10^{-15}$~Tg
    $\approx$ 8~250~Tg/a\~no]}

\end{ejercicios}

\end{bloque_ejercicios}
% =============================================================================
% FIN DE ejercicios_cap09.tex
% =============================================================================
